% ============ 第1章 §1.11 无穷数列情形的推广 (原书 pp.33-35) ============

\section{无穷数列情形的推广}\label{sec:1.11}
\index{J@基本向量}




设 $N_{0} = \{0, 1, \ldots\}$. n 再次设为自然数. 这次考虑数列项数不是 $2^{nN}$ 个而是无限多个的情形. 即, 考虑表示为 $\{a_{\vec{j}}\}_{\vec{j}\in\mathbb{N}_{0}^{n}}$ 形式的多重数列. 对 $j \in \mathbb{N}$, $\mathbf{e}_{j} = \{\chi_{\{j\}}(k)\}_{k\in\mathbb{N}_{0}}$ 表示基本向量. 将 $E_{k}$ 视为取 $I(k,j)$ ($j = 0, 1, \ldots$) 上的平均, 与此对应, $\vec{X}_{k}, k \in \mathbb{N}$ 也自然地被推广. 对 $1 < p < \infty$, 由于迄今定理中出现的常数 $C_{p}$ 不依赖于 N, 令
\[
g(a)_{\vec{j}} = \left(\sum_{k=1}^{\infty}\sum_{\vec{X}\in\{(D,E)\}^{n}\setminus\{\mathbf{E}\}}|\vec{X}_{k-1}(a)_{\vec{j}}|^{2}\right)^{1/2} \quad (\vec{j} \in \mathbb{N}_{0}^{n}).
\]
对 $1 \leq p \leq \infty$, 多重数列空间 $\ell^{p}(\mathbb{N}_{0}^{n})$ 指具有 $\{a_{\vec{j}}\}_{\vec{j}\in\mathbb{N}_{0}^{n}}$ 形式且满足
\[
\|\{a_{\vec{j}}\}_{\vec{j}\in\mathbb{N}_{0}^{n}}\|_{\ell^{p}} = \left(\sum_{\vec{j}\in\mathbb{N}_{0}^{n}}|a_{\vec{j}}|^{p}\right)^{1/p} < \infty
\]
的数列全体. 当 $p = \infty$ 时, 右边不取 p 次幂和而取上限. 当 n = 1 时, 记为 $\ell^{p}(\mathbb{N}_{0})$. 将指标集合 $N_{0}$ 换成 $\mathbb{N}$ 所得的数列空间 $\ell^{p}(\mathbb{N})$ 也同样定义.

\begin{lemma}\label{lem:1.5}
对 $a = \{a_{\vec{j}}\}_{\vec{j}\in\mathbb{N}_{0}^{n}} \in \ell^{2}(\mathbb{N}_{0}^{n})$, 成立
\[
\|a\|_{\ell^{2}}^{2} = \sum_{k=1}^{\infty}\sum_{\vec{X}\in\{(D,E)\}^{n}\setminus\{\mathbf{E}\}}\|\vec{X}_{k-1}(a)\|_{\ell^{2}}^{2}.
\]
\end{lemma}

\begin{proof}
固定自然数 $N_{0}$. 由 $(A+B)^{2} \leq 2A^{2} + 2B^{2}$, $A, B \in \mathbb{R}$, 与 Cauchy--Schwarz 不等式, 若自然数 N 满足 $N > N_{0}$, 则利用式~\eqref{eq:1.14} 给出的区间 $I(N,0), I(N_{0},0)$, 有
\begin{align*}
\frac{1}{2^{nN}}\left|\sum_{\vec{j}\in I(N,0)^{n}}a_{\vec{j}}\right|^{2}
&\leq \frac{2}{2^{nN}}\left|\sum_{\vec{j}\in I(N_{0},0)^{n}}a_{\vec{j}}\right|^{2} + \frac{2}{2^{nN}}\left|\sum_{\vec{j}\in I(N,0)^{n}\setminus I(N_{0},0)^{n}}a_{\vec{j}}\right|^{2} \\
&\leq \frac{2}{2^{nN}}\left|\sum_{\vec{j}\in I(N_{0},0)^{n}}a_{\vec{j}}\right|^{2} + 2\sum_{\vec{j}\in I(N,0)^{n}\setminus I(N_{0},0)^{n}}|a_{\vec{j}}|^{2}.
\end{align*}
对 N 取上极限, 得
\[
\limsup_{N\to\infty}\frac{1}{2^{nN}}\left|\sum_{\vec{j}\in I(N,0)^{n}}a_{\vec{j}}\right|^{2} \leq 2\sum_{\vec{j}\in\mathbb{N}_{0}^{n}\setminus I(N_{0},0)^{n}}|a_{\vec{j}}|^{2} = 2\sum_{\vec{j}\in\mathbb{N}_{0}^{n}}|a_{\vec{j}}|^{2} - 2\sum_{\vec{j}\in I(N_{0},0)^{n}}|a_{\vec{j}}|^{2}.
\]
再对 $N_{0}$ 取极限, 得
\[
\limsup_{N\to\infty}\frac{1}{2^{nN}}\left|\sum_{\vec{j}\in I(N,0)^{n}}a_{\vec{j}}\right|^{2} = 0 \quad \text{即} \quad \lim_{N\to\infty}\frac{1}{2^{nN}}\left|\sum_{\vec{j}\in I(N,0)^{n}}a_{\vec{j}}\right|^{2} = 0.
\]
这里, 令 $a^{N} = \{\chi_{I(N,0)^{n}}(\vec{j})a_{\vec{j}}\}_{\vec{j}\in\mathbb{N}_{0}^{n}}$, 则由定理~\ref{thm:1.9} 的 n 维版本,
\[
\|a^{N}\|_{\ell^{2}}^{2} = \frac{1}{2^{nN}}\left|\sum_{\vec{j}\in I(N,0)^{n}}a_{\vec{j}}\right|^{2} + \sum_{k=1}^{N}\sum_{\vec{X}\in\{(D,E)\}^{n}\setminus\{\mathbf{E}\}}\|\vec{X}_{k-1}(a)\|_{\ell^{2}}^{2},
\]
令 $N \to \infty$ 即得结论.
\end{proof}

对 $p \in (1, \infty)\setminus\{2\}$ 也仿照刚才的讨论, 在定理~\ref{thm:1.14} 中令 $N \to \infty$, 得下面的定理.

\begin{theorem}\label{thm:1.15}
存在仅依赖于自然数 n 与 $p \in (1, \infty)$ 的常数 $C > 0$, 使得对所有 $a \in \ell^{p}(\mathbb{N}_{0}^{n})$ 成立
\[
C^{-1}\|a\|_{\ell^{p}} \leq \|g(a)\|_{\ell^{p}} \leq C\|a\|_{\ell^{p}}.
\]
\end{theorem}

把``存在常数 $C > 0$, 使得对所有 $a \in \ell^{p}(\mathbb{N}_{0}^{n})$ 成立 $C^{-1}\|a\|_{\ell^{p}} \leq \|g(a)\|_{\ell^{p}} \leq C\|a\|_{\ell^{p}}$''这一状况记为 $\|g(a)\|_{\ell^{p}} \sim \|a\|_{\ell^{p}}$ 是很方便的. 此后, 在类似的情形下也常用 $\sim$ 表示. 又, 当只考虑 $\|g(a)\|_{\ell^{p}} \leq C\|a\|_{\ell^{p}}$ 这种单侧情形时, 记为 $\|g(a)\|_{\ell^{p}} \lesssim \|a\|_{\ell^{p}}$ 或 $\|a\|_{\ell^{p}} \gtrsim \|g(a)\|_{\ell^{p}}$.

\begin{exercise}\label{exer:1.28}
设数列 $\{a_{j}\}_{j=1}^{\infty} \in \ell^{1}(\mathbb{N})$ 对所有 $j \in \mathbb{N}$ 满足 $0 \leq a_{j+1} \leq a_{j}$. 此时, 证明 $\lim_{j\to\infty}ja_{j} = 0$.
\end{exercise}

\begin{exercise}\label{exer:1.29}
设 $1 \leq p \leq \infty$.

(1) 当 $0 \leq j < k < \infty$ 时, 计算 $\|\mathbf{e}_{j} - \mathbf{e}_{k}\|_{\ell^{p}}$.

(2) 证明 $\ell^{p}(\mathbb{N})$ 的单位球 $B = \{a \in \ell^{p}(\mathbb{N}) : \|a\|_{\ell^{p}} \leq 1\}$ 不是紧的.
\end{exercise}

\begin{exercise}\label{exer:1.30}
设 $1 \leq p \leq \infty$. 用 $\operatorname{Span}(\{e_{j}\}_{j=0}^{\infty})$ 表示包含各 $e_{j}$ 的最小线性空间. 此时, 对配备 $\ell^{\infty}$ 范数的 $\operatorname{Span}(\{e_{j}\}_{j=0}^{\infty})$, 证明以下.

(1) 若 $x_{k} = \sum_{l=1}^{k}2^{-l}\mathbf{e}_{l}$, 则当 $1 \leq k < k' < \infty$ 时, $\|x_{k} - x_{k'}\|_{\ell^{\infty}} = 2^{-k-1}$.

(2) 在 $\operatorname{Span}(\{\mathbf{e}_{j}\}_{j=0}^{\infty})$ 中, $\{x_{k}\}_{k=1}^{\infty}$ 不收敛.
\end{exercise}
